% -------------------------------------------------------
%  AQI-MLops – Conference Paper (ACM sigconf style)
%  Author: Daniel Lai, Drexel University, May 2026
% -------------------------------------------------------
\PassOptionsToPackage{hyphens}{url}
\documentclass[sigconf, nonacm]{acmart}

% ---- packages -----------------------------------------
\usepackage{booktabs}
\usepackage{array}
\usepackage{amsmath}
\usepackage{microtype}
\usepackage{listings}
\usepackage{xcolor}
\usepackage{url}

% ---- allow long paths / URLs to break everywhere ------
\urlstyle{tt}
\def\UrlBreaks{\do\/\do\-\do\_\do\.\do\:\do\=\do\?}

% ---- give LaTeX extra stretch to avoid overfull hboxes
\emergencystretch=4em
\tolerance=1500
\hbadness=1500

% ---- listing style ------------------------------------
\lstset{
  basicstyle=\ttfamily\small,
  breaklines=true,
  breakatwhitespace=false,
  frame=single,
  backgroundcolor=\color{gray!8},
  columns=flexible,
  keepspaces=true,
}

% ---- hyperref (acmart loads it; just set colours) -----
\AtBeginDocument{%
  \hypersetup{colorlinks=true, urlcolor=blue,
              linkcolor=black, citecolor=black}}

% =======================================================
\begin{document}

\title{AQI-MLops: An End-to-End Production MLOps Pipeline
       for Real-Time Air Quality Index Prediction
       Across 99 Global Cities}

\author{Daniel Lai}
\affiliation{%
  \institution{Drexel University, College of Computing \& Informatics}
  \city{Philadelphia}
  \state{PA}
  \country{USA}}
\email{dl3528@drexel.edu}

% =======================================================
\begin{abstract}
Air quality affects the health of billions of people, yet
accessible real-time forecasts at global scale remain scarce.
This paper presents \textbf{AQI-MLops}, a fully automated,
cloud-native machine learning operations (MLOps) system that
collects live air quality sensor readings every hour for
99 cities worldwide, stores them in a medallion data lakehouse
on AWS, predicts the next-hour Air Quality Index (AQI) class
using a tuned XGBoost classifier, and serves those predictions
through a web dashboard with sub-millisecond response times.
The system achieves a cross-validated macro F1 score of
\textbf{0.9425} and a validation macro F1 score of
\textbf{0.9535} across five AQI classes.  The full pipeline,
from raw API ingest to model retraining and observability,
runs continuously and autonomously using AWS Lambda, S3,
Athena (Apache Iceberg), Redis, FastAPI, Docker, and Grafana
Cloud, with CI/CD orchestrated by GitHub Actions.
\end{abstract}

\keywords{MLOps, Air Quality Index, XGBoost, AWS,
          Apache Iceberg, FastAPI, Observability}

\maketitle

% =======================================================
\section{Introduction}

The World Health Organization estimates that 99\% of the
global population breathes air that exceeds recommended
quality limits~\cite{who2021}.  Timely, city-level AQI
forecasts can drive behavioral decisions: whether to
exercise outdoors, whether to issue public health alerts,
or whether to activate industrial emission controls.
Despite the importance of this information, most publicly
available AQI data is either delayed batch reports or
single-city solutions that do not generalize.

This work addresses the challenge by building a
production-grade MLOps system with three primary properties:

\begin{enumerate}
  \item \textbf{Scale}: 99 geographically diverse cities
        spanning Asia, Europe, the Americas, and Africa.
  \item \textbf{Freshness}: Hourly ingest, hourly prediction
        refresh, and automated daily model retraining.
  \item \textbf{Reliability}: ACID-consistent data storage,
        Redis-backed sub-millisecond serving, self-healing
        cache warming, and a full observability stack.
\end{enumerate}

The system is implemented in Python and deployed on AWS, with
all infrastructure described as configuration rather than
hand-operated scripts.  All code is publicly available at
\url{https://github.com/danhdanhtuan0308/AQI-MLops}.

% =======================================================
\section{Related Work}

Traditional air quality forecasting relies on numerical
weather prediction models such as CMAQ and WRF-Chem, which
are computationally expensive and require domain expertise
to configure~\cite{cmaq2006}.  Data-driven approaches using
gradient-boosted trees and recurrent neural networks have
demonstrated competitive accuracy on city-level
datasets~\cite{zhao2018,qi2019}, but published systems
rarely describe the end-to-end engineering infrastructure
required to keep such models fresh and serving at production
quality.

MLOps as a discipline formalizes the practices needed to
operationalize ML models~\cite{sculley2015}.  Systems such
as Uber Michelangelo~\cite{michelangelo2017} and Meta
FBLearner Flow~\cite{fblearner2016} demonstrate
industrial-scale MLOps but are not open-source and operate
at a scale impractical for academic projects.  AQI-MLops
occupies a middle ground: a real production system with real
traffic, open-source code, and infrastructure costs within
a research budget.

% =======================================================
\section{System Architecture}

The system is composed of five sequential pipeline stages
executed across three compute substrates: AWS Lambda, AWS
EC2, and GitHub Actions.
Figure~\ref{fig:pipeline} illustrates the full data flow.

% --- full-width architecture diagram -------------------
\begin{figure*}[t]
\small
\begin{verbatim}
OpenWeatherMap API
  |
  +-- fetch_weather.py  (bulk backfill, 1 year x 99 cities)
  |         |
  |         +-> S3 Bronze: raw_pipeline  (Parquet / SNAPPY)
  |
  +-- Lambda A  (hourly EventBridge cron)
            |
            +-> S3 Bronze: raw_hourly  (Parquet, Hive-partitioned)
                      |
                  Lambda B: MERGE INTO aqi_unified  (Iceberg, ACID)
                      |
                  Silver: aqi_db.aqi_unified
                      |
                  FastAPI /warm-cache
                      |
                  XGBoost inference (19 features) --> Redis  (TTL 2 h)
                      |
                  Lambda C  (hourly + 5 min offset)
                      |
                  Gold: predictions Iceberg table
                      |
                  GitHub Actions  (daily cron)
                      |
                  ml/train.py --> MLflow --> model-registry/model.ubj
\end{verbatim}
\caption{End-to-end pipeline data flow across five sequential stages.}
\label{fig:pipeline}
\end{figure*}

% -------------------------------------------------------
\subsection{Data Ingestion}

\textbf{Batch Backfill.}
The script \path{data-pipeline/fetch_weather.py} seeds the
lakehouse with one year of historical data
(2025-03-20 to 2026-03-20) for all 99 cities.  The date
range is split into seven-day chunks per city to respect
the OpenWeatherMap (OWM) free-tier pagination limit.  A
0.5-second inter-request sleep bounds throughput to
approximately 2 requests per second, comfortably within
OWM free-tier rate limits.  Each chunk is serialized as
SNAPPY-compressed Parquet via PyArrow and uploaded to
\path{s3://weather-bulk/data-pipeline/} via the AWS SDK.

\textbf{Live Hourly Ingest (Lambda A).}
An AWS EventBridge rule fires Lambda A every hour.
Lambda A calls the OWM Current Air Pollution API for all
99 cities and writes a single Parquet file to a
Hive-partitioned path of the form
\path{s3://weather-bulk/hourly/year=YYYY/month=MM/day=DD/}.
On success, Lambda A invokes Lambda B asynchronously,
passing the S3 key of the new file.

The raw data schema captures eleven air quality measurements
alongside city metadata, as shown in Table~\ref{tab:schema}.

\begin{table}[h]
  \caption{Raw data schema (key columns).}
  \label{tab:schema}
  \small
  \begin{tabular}{@{}p{2.0cm}p{5.2cm}@{}}
    \toprule
    \textbf{Column} & \textbf{Type and Description} \\
    \midrule
    \texttt{timestamp}  & TIMESTAMP (UTC): observation time \\[2pt]
    \texttt{city\_slug} & STRING: URL-safe city identifier \\[2pt]
    \texttt{aqi}        & INT (1 to 5): overall AQI class \\[2pt]
    \texttt{co}, \texttt{no2}, \texttt{o3},
    \texttt{so2}, \texttt{pm2\_5},
    \texttt{pm10}, \texttt{nh3}
      & DOUBLE ($\mu$g/m$^{3}$): pollutant concentrations \\
    \bottomrule
  \end{tabular}
\end{table}

% -------------------------------------------------------
\subsection{Data Lakehouse: Medallion Architecture}

The storage layer is built on AWS S3 with Apache Iceberg
tables registered in AWS Glue Data Catalog and queried via
AWS Athena.

\begin{itemize}
  \item \textbf{Bronze Layer.}  Two read-only external
        tables: \texttt{aqi\_db.raw\_pipeline} (bulk batch)
        and \texttt{aqi\_db.raw\_hourly} (live hourly,
        Hive-partitioned by year, month, and day).

  \item \textbf{Silver Layer.}  Lambda B executes an Athena
        \texttt{MERGE INTO aqi\_db.aqi\_unified} statement
        keyed on \texttt{(city\_slug, timestamp)}.  This
        provides ACID-safe deduplication and ensures
        idempotency across overlapping pipeline runs.
        A \texttt{source} column (value \texttt{pipeline}
        or \texttt{live}) is written for data lineage.
        The Iceberg format enables time-travel queries and
        schema evolution without table rewrites.

  \item \textbf{Gold Layer.}  Aggregated analytics tables
        stored at
        \path{s3://weather-bulk/process/BI/},
        used by Grafana for business-intelligence dashboards.
\end{itemize}

% -------------------------------------------------------
\subsection{Inference and Serving}

When Lambda B completes the merge, it sends a
\texttt{POST /warm-cache} request to the FastAPI service
running on EC2.  On receiving this call, FastAPI:

\begin{enumerate}
  \item Queries Athena for the last 340 hours of sensor
        data across all 99 cities in a single batch query.
  \item Builds 19-dimensional feature vectors for each city
        (see Section~\ref{sec:features}).
  \item Runs batch XGBoost inference.
  \item Writes all 99 city predictions to Redis with a
        two-hour TTL.
  \item Persists predictions to \texttt{aqi:pred\_buffer}
        (a Redis list) for Lambda C to flush to Iceberg.
\end{enumerate}

User-facing \texttt{GET /predict/\{city\_slug\}} requests
are served entirely from Redis, so no Athena query occurs
at request time, yielding sub-millisecond response latency.
A background thread calls \texttt{warm\_cache()} every
3\,600 seconds independently of Lambda B, ensuring cache
freshness even if Lambda B is unavailable.

% -------------------------------------------------------
\subsection{Prediction Persistence (Lambda C)}

Lambda C runs at five minutes past every hour via
EventBridge.  It reads the \texttt{aqi:pred\_buffer} list
from Redis, writes the records as Parquet to
\path{s3://weather-bulk/predictions/},
and runs \texttt{INSERT INTO} on a dedicated Iceberg
\texttt{predictions} table.  This builds a permanent,
queryable record of every prediction the model has ever
made, enabling retrospective drift analysis.

% -------------------------------------------------------
\subsection{Model Retraining and CI/CD}

A GitHub Actions workflow (\path{daily_retrain.yml}) runs
every day.  It invokes \path{ml/train.py}, which reads
hyperparameters from \path{ml/best-config.yaml}, trains
XGBoost on the full Athena dataset, logs all parameters
and metrics to MLflow, and writes the new model artifact
to \path{ml/model-registry/model.ubj}.

The FastAPI endpoint \texttt{GET /reload-model}
hot-swaps the in-memory model without restarting the
container.  On every merge to \texttt{main},
\path{cd.yml} rebuilds the Docker image and redeploys
to EC2 via \path{deploy/redeploy.sh}.

% =======================================================
\section{Data Science}

% -------------------------------------------------------
\subsection{Exploratory Data Analysis}

The AQI class distribution across the training corpus is
shown in Table~\ref{tab:classdist}.  Classes 2 and 3
dominate, reflecting that most cities spend most hours in
mild-to-moderate conditions.  Classes 4 and 5, while
minority, are the most safety-relevant and were addressed
through class-weighted training.

\begin{table}[h]
  \caption{AQI class distribution in training data.}
  \label{tab:classdist}
  \small
  \begin{tabular}{@{}clr@{}}
    \toprule
    \textbf{Class} & \textbf{Label} & \textbf{Proportion} \\
    \midrule
    1 & Good      & 11.8\% \\
    2 & Fair      & 38.7\% \\
    3 & Moderate  & 22.6\% \\
    4 & Poor      & 14.0\% \\
    5 & Very Poor & 12.9\% \\
    \bottomrule
  \end{tabular}
\end{table}

\textbf{Data Quality Issues.}
OWM occasionally returns AQI class 6 for extreme pollution
events outside their documented scale; these labels are
clamped to 5 during data loading.  Lag features require at
least three prior hours of data per city, so the first three
rows of each city's time series are dropped before training.
Missing mid-series hours break the lag chain and are
also dropped.

% -------------------------------------------------------
\subsection{Feature Engineering}
\label{sec:features}

AQI exhibits strong temporal autocorrelation: the next
hour's AQI class is highly predictable from recent history.
The feature set (19 features total) is organized into three
groups.

\textbf{Momentum Features} capture temporal autocorrelation
and are listed in Table~\ref{tab:features}.

\begin{table}[h]
  \caption{Momentum features (7 of 19 total).}
  \label{tab:features}
  \small
  \begin{tabular}{@{}p{1.9cm}p{2.6cm}p{2.6cm}@{}}
    \toprule
    \textbf{Feature} & \textbf{Formula} & \textbf{Semantic} \\
    \midrule
    \texttt{pm10\_lag1}
      & PM10 at $T{-}1$
      & Prior particulate \\[2pt]
    \texttt{aqi\_delta\_1h}
      & $\mathrm{AQI}(T){-}\mathrm{AQI}(T{-}1)$
      & 1 h AQI trend \\[2pt]
    \texttt{aqi\_delta\_2h}
      & $\mathrm{AQI}(T){-}\mathrm{AQI}(T{-}2)$
      & 2 h AQI momentum \\[2pt]
    \texttt{aqi\_delta\_3h}
      & $\mathrm{AQI}(T){-}\mathrm{AQI}(T{-}3)$
      & Sustained trend \\[2pt]
    \texttt{pm10\_delta\_1h}
      & $\mathrm{PM10}(T){-}\mathrm{PM10}(T{-}1)$
      & Particulate change \\[2pt]
    \texttt{pm25\_delta\_1h}
      & $\mathrm{PM2.5}(T){-}\mathrm{PM2.5}(T{-}1)$
      & Fine part. change \\[2pt]
    \texttt{aqi\_roll\_std4}
      & $\sigma(\mathrm{AQI}_{T-3..T})$
      & 4 h volatility \\
    \bottomrule
  \end{tabular}
\end{table}

\textbf{Cyclical Time Encoding.}
Raw hour (0 to 23) and month (1 to 12) integers break at
natural boundaries: hour 23 and hour 0 are adjacent in time
but numerically distant.  Sine and cosine encoding resolves
this:
\begin{align}
  \mathrm{hour\_sin} &= \sin\!\left(\tfrac{2\pi h}{24}\right),
  &\mathrm{hour\_cos} &= \cos\!\left(\tfrac{2\pi h}{24}\right)\\
  \mathrm{month\_sin} &= \sin\!\left(\tfrac{2\pi m}{12}\right),
  &\mathrm{month\_cos} &= \cos\!\left(\tfrac{2\pi m}{12}\right)
\end{align}

\textbf{Point-in-Time Pollutants.}
The remaining eight features are the current-hour
concentrations of \texttt{co}, \texttt{no}, \texttt{no2},
\texttt{o3}, \texttt{so2}, \texttt{nh3}, \texttt{pm10},
and \texttt{pm25\_ratio} (PM2.5 normalized by total
pollutant mass).  Missing values are filled with
per-feature medians persisted to
\path{ml/model-registry/median.json}.

% -------------------------------------------------------
\subsection{Model Selection and Hyperparameter Tuning}

Three model families were benchmarked: Logistic Regression
(baseline), CatBoost, and XGBoost.  XGBoost consistently
achieved the highest macro F1 across cross-validation folds
and was selected for production.

Hyperparameter optimization used Random Search (20 trials,
3-fold time-series cross-validation) and Bayesian Search
(20 trials).  Results are shown in Table~\ref{tab:hpo}.

\begin{table}[h]
  \caption{Hyperparameter optimization comparison.}
  \label{tab:hpo}
  \small
  \begin{tabular}{@{}p{2.6cm}rr@{}}
    \toprule
    \textbf{Method} & \textbf{Best CV F1} & \textbf{Trials} \\
    \midrule
    Random Search   & \textbf{0.9425} & 20 \\
    Bayesian Search & 0.9410          & 20 \\
    \bottomrule
  \end{tabular}
\end{table}

Random Search was selected as the production strategy.
The best hyperparameters are shown in
Table~\ref{tab:bestparams}.

\begin{table}[h]
  \caption{Best XGBoost hyperparameters (Random Search).}
  \label{tab:bestparams}
  \small
  \begin{tabular}{@{}p{2.2cm}rp{2.2cm}r@{}}
    \toprule
    \textbf{Parameter} & \textbf{Value}
      & \textbf{Parameter} & \textbf{Value} \\
    \midrule
    colsample\_bytree & 0.5997 & learning\_rate     & 0.0810 \\
    gamma             & 0.1516 & max\_delta\_step   & 1      \\
    max\_depth        & 6      & min\_child\_weight & 8      \\
    n\_estimators     & 542    & reg\_alpha         & 0.9359 \\
    reg\_lambda       & 6.4496 & subsample          & 0.6599 \\
    \bottomrule
  \end{tabular}
\end{table}

\textbf{Class-Weighted Training.}
To handle class imbalance and improve prediction of
safety-critical classes 4 and 5, a custom sample weight
scheme was applied:
\begin{itemize}
  \item Sustained AQI transitions (new class held for two
        or more consecutive hours): $5\times$ weight boost.
  \item Single-hour AQI blips (new class reverts next
        hour): $2\times$ weight boost.
\end{itemize}
Validation macro F1 is \textbf{0.9535}.  All training runs
are tracked in MLflow under the experiment
\texttt{AQI-Classification}, with the best model registered
in the \texttt{AQI-XGBoost} model registry.

% =======================================================
\section{MLOps and Production Engineering}

% -------------------------------------------------------
\subsection{Containerization and Deployment}

The FastAPI service is packaged as a Docker image built
from the project \texttt{Dockerfile} and managed by
\texttt{docker-compose.yml}.  Application source and model
artifacts are mounted as volumes, decoupling code updates
from image rebuilds.  Nginx proxies port 80 to the internal
FastAPI port 8000.  A \texttt{systemd} unit file
(\path{deploy/aqi-api.service}) ensures the container
restarts automatically on crash or machine reboot.

% -------------------------------------------------------
\subsection{Observability}

The system pushes all telemetry signals directly to Grafana
Cloud without a self-hosted collection layer, using Basic
auth via an OTLP gateway:

\begin{itemize}
  \item \textbf{Metrics (Grafana Mimir)} via the
        OpenTelemetry SDK: prediction request count,
        latency histogram, cache hit/miss rate, model drift
        scores per city, and system CPU/memory.

  \item \textbf{Logs (Grafana Loki)} via an async queue
        handler: structured JSON logs from all service
        components.

  \item \textbf{Traces (Grafana Tempo)} via OTLP HTTP
        exporter: distributed request traces from
        \texttt{FastAPIInstrumentor}, with excluded URLs
        for health and metrics endpoints to reduce noise.
\end{itemize}

Alerts are configured with threshold-based rules that push
notifications to Teams, Slack, or Discord when drift or
error rates exceed defined bounds.

% -------------------------------------------------------
\subsection{Testing and CI}

The test suite under \path{tests/} covers three layers:

\begin{itemize}
  \item \path{test_api.py}: HTTP endpoint tests via the
        FastAPI \texttt{TestClient}.
  \item \path{test_inference.py}: feature pipeline and
        model inference correctness.
  \item \path{test_training_logic.py}: training logic
        validation including class weight computation and
        feature array shapes.
\end{itemize}

GitHub Actions \path{ci.yml} runs linting and the full
test suite on every pull request.  Merges to \texttt{main}
trigger \path{cd.yml}, which rebuilds and redeploys the
service to EC2.

% =======================================================
\section{Results}

\begin{table}[h]
  \caption{Production system performance summary.}
  \label{tab:results}
  \small
  \begin{tabular}{@{}p{3.7cm}p{3.5cm}@{}}
    \toprule
    \textbf{Metric} & \textbf{Value} \\
    \midrule
    Cross-validated macro F1    & 0.9425 \\
    Validation macro F1         & 0.9535 \\
    Cities covered              & 99 \\
    Features at inference       & 19 \\
    Prediction latency (p50)    & $<$1 ms (Redis) \\
    Warm-cache query latency    & 3 to 8 seconds \\
    Cache TTL                   & 2 hours \\
    Retraining frequency        & Daily (GitHub Actions) \\
    MLflow model versions       & 11 or more \\
    \bottomrule
  \end{tabular}
\end{table}

The dominant error cases are transitions between adjacent
classes, which is expected given their similar pollutant
concentrations.  The class-weighting scheme substantially
improved recall for classes 4 and 5 relative to the
unweighted baseline.

% =======================================================
\section{Discussion}

\textbf{Design Decisions.}
Pre-computing and caching all 99 city predictions at ingest
time, rather than at request time, was the single most
impactful engineering decision for serving quality.  It
shifts the latency cost to the data pipeline layer where it
is acceptable, and makes user-facing latency independent of
model complexity or Athena query time.

\textbf{Limitations.}
The current feature set does not include wind direction or
speed, which are known drivers of pollutant
dispersion~\cite{cmaq2006}.  The training corpus covers a
single year (2025 to 2026); longer histories would improve
seasonal generalization.  Relying on a single upstream API
means undocumented API behaviors propagate silently unless
caught at ingestion.

\textbf{Future Work.}
Planned extensions include: (1)~incorporating
meteorological features such as wind and precipitation to
improve accuracy for classes 4 and 5; (2)~adding a
geographic embedding to capture inter-city correlation;
(3)~replacing the Redis TTL-based cache with a streaming
architecture (Kafka) to push predictions to clients in real
time; and (4)~implementing automated A/B testing between
model versions in the registry.

% =======================================================
\section{Conclusion}

AQI-MLops demonstrates that a full production MLOps system
covering data engineering, feature engineering, model
training, experiment tracking, containerized deployment,
Redis-backed serving, and cloud-native observability can be
built and maintained by a small team as an academic capstone
project.  The system achieves 94\%+ macro F1 on a
five-class air quality classification task across 99 global
cities, with fully automated daily retraining and
sub-millisecond prediction serving.  All code, pipeline
definitions, model configurations, and infrastructure
scripts are publicly available at the repository below.

% =======================================================
\section*{Repository}

\textbf{GitHub:}
\url{https://github.com/danhdanhtuan0308/AQI-MLops}

\medskip
\noindent
The repository is organized as follows:

\begin{lstlisting}[basicstyle=\footnotesize\ttfamily,
                   frame=none, backgroundcolor=\color{white}]
AQI-MLops/
  app/             FastAPI service and telemetry
    main.py          Routes, cache warming, background thread
    inference.py     Feature building, XGBoost predict
    telemetry.py     OTel tracing, Mimir, Loki logging
  data-pipeline/   Batch backfill and live ingest
    fetch_weather.py OWM API to S3 Parquet
    lambda/
      handler.py     Lambda A hourly live ingest
  lakehouse/       DDL and merge logic
    setup.sql        Athena DDL: Iceberg tables
    merge_handler.py Lambda B: MERGE INTO aqi_unified
    prediction_flush_handler.py  Lambda C Redis flush
  ml/              Training pipeline and model registry
    train.py         Feature eng, XGBoost, MLflow logging
    best-config.yaml Best hyperparameters
    model-registry/  model.ubj, features.json, median.json
  experiement/     EDA and benchmarking notebooks
  tests/           Unit and integration tests
  deploy/          EC2 setup, Nginx, systemd service
  .github/         CI, CD, and daily retraining workflows
  Dockerfile
  docker-compose.yml
  pyproject.toml
\end{lstlisting}

% =======================================================
\begin{thebibliography}{7}

\bibitem{who2021}
World Health Organization.
\textit{WHO Global Air Quality Guidelines.}
Geneva: WHO, 2021.

\bibitem{cmaq2006}
Byun, D., \& Schere, K.~L. (2006).
Review of the governing equations, computational algorithms,
and other components of the Models-3 Community Multiscale
Air Quality (CMAQ) modeling system.
\textit{Applied Mechanics Reviews}, 59(2), 51--77.

\bibitem{zhao2018}
Zhao, R., Gu, X., Xue, B., Zhang, J., \& Ren, W. (2018).
Short period PM2.5 prediction based on multivariate linear
regression model.
\textit{PLOS ONE}, 13(7), e0201011.

\bibitem{qi2019}
Qi, Z., Wang, T., Song, G., Hu, W., Li, X., \& Zhang, Z.
(2019).
Deep air learning: Interpolation, prediction, and feature
analysis of fine-grained air quality.
\textit{IEEE Transactions on Knowledge and Data
Engineering}, 32(12), 2285--2297.

\bibitem{sculley2015}
Sculley, D., et al. (2015).
Hidden technical debt in machine learning systems.
\textit{Advances in Neural Information Processing
Systems}, 28.

\bibitem{michelangelo2017}
Hermann, J., \& Del Balso, M. (2017).
Meet Michelangelo: Uber's Machine Learning Platform.
\textit{Uber Engineering Blog}.

\bibitem{fblearner2016}
Dunn, J. (2016).
Introducing FBLearner Flow: Facebook's AI backbone.
\textit{Facebook Engineering Blog}.

\end{thebibliography}

\end{document}
